\subsection{System Architecture}
\label{subsec:methodology-system-overview}

This work presents a control framework for quadruped robots capable
of generating dynamic, acyclic gaits in challenging environments
without requiring a predefined periodic schedule. The framework uses
a hierarchical architecture that separates high-level footstep
decision-making from low-level kinematic tracking. At each control
tick, a terrain height scan is thresholded and dilated into a
per-leg validity mask, a small set of candidate footsteps is drawn
within that mask and locally refined against GaitNet's learned value
function, and the single best candidate\textemdash or a no-action
option\textemdash is selected for the entire robot and passed to a
low-level tracking controller for execution via model predictive
control and whole-body feedback. GaitNet's footstep action is
consumed identically regardless of which controller executes it:
in simulation, the tracking controller is the MIT Mini-Cheetah
Controller \cite{di_mini_cheetah_2020}, structured in python by
\cite{zhuang_rl_mpc_locomotion_2025}; on hardware, GaitNet instead
supplies footstep targets directly to the Legged Perceptive Stack's
own NMPC/Whole-Body Control (WBC) pipeline \cite{leggedcontrol}, in
place of that framework's native gait scheduler.

Two constraints are enforced directly in the validity mask, before
GaitNet ever sees a candidate. First, a leg already in its swing
phase is excluded from candidate generation, since it cannot be
redirected mid-swing. Second, a new step is permitted only while at
least \numberstringnum{\corlMinContactLegs} legs remain in ground
contact, which caps concurrent swing at
\numberstringnum{\corlMaxConcurrentSwingLegs} legs; this is a
deliberate feasibility constraint, not an incidental limitation, and
every result in \autoref{sec:results} operates under it. The validity
mask itself is computed from a thresholded height scan, dilated by
\numberstringnum{\corlMaskDilationRounds} rounds of max-pooling to
conservatively account for the physical extent of the robot's foot.

\begin{figure*}[t]
  \centering
  \includegraphics[width=0.7\linewidth]{images/diagrams/control-system.png}
  \caption{System architecture of the proposed control framework. A
    terrain height scan is thresholded and dilated into a per-leg
    validity mask; a small set of candidate footsteps is sampled
    within that mask per leg and locally refined against GaitNet's
    value function; the single best candidate across all legs, or a
    no-action option, is selected and passed to the tracking
  controller.}
  \label{fig:diagram-control-system}
\end{figure*}

\subsection{High-Rate Action Selection}
\label{subsec:methodology-high-rate-selection}

GaitNet is evaluated at \corlControlHz\,Hz, a \corlControlTickMs\,ms
control tick, and never represents multi-leg coordination directly:
at most one footstep is committed per tick, for one leg, and the rest
of the robot is left untouched until its own turn. A selected swing
lasts between \corlSwingDurationMin\,s and \corlSwingDurationMax\,s,
so \numberstringnum{\corlMinTicksPerSwing} to
\numberstringnum{\corlMaxTicksPerSwing} further decision ticks elapse
while that leg is still airborne. GaitNet may therefore commit a
second leg
mid-swing of the first, subject only to the contact-count constraint
above. Multi-leg concurrency is an emergent consequence of
re-deciding faster than a swing completes, not a behavior the
architecture was given machinery to plan; whether the policy learns
to anticipate future ticks and implicitly coordinate legs through
this repeated re-selection is a question of what reinforcement
learning discovers, not something encoded in the action space itself.
This is the central mechanism of the system: the simplicity of a
single-leg, single-tick action space is what makes acyclic, dynamic
gait generation tractable at control-loop rates.

\subsection{GaitNet Formulation}
\label{subsec:methodology-gaitnet-architecture}

GaitNet is a single-network architecture comprised of a robot-state
encoder, a footstep encoder, a shared trunk, and two task-specific
output heads. The first head predicts a scalar value
logit used to rank a candidate; the second predicts the swing
duration to use if that candidate is selected. GaitNet evaluates
candidates conditioned on a one-hot leg encoding (including a
no-action option) and the current tracking state $\mathbf{x_g}$:

\begin{equation}
  \mathbf{x_g} =
  \begin{bmatrix}
    \mathbf p_{b,xy} &
    \mathbf r_{w,z} &
    \mathbf v_b &
    \mathbf \omega_b &
    \mathbf u &
    \mathbf g &
    \mathbf c
  \end{bmatrix}^T
  \label{eq:gaitnet-x}
\end{equation}

Here $\mathbf p_{b,xy} \in \mathbb{R}^{\corlFootPosDim}$ is the $xy$
position of all four end-effectors in the base frame, and
$\mathbf r_{w,z}$ is the height of the robot's center of mass in the
world frame. $\mathbf v_b \in \mathbb{R}^{\corlVelDim}$ and
$\mathbf \omega_b \in \mathbb{R}^{\corlAngVelDim}$ are the robot's
linear and angular velocity in the root frame.
$\mathbf u = [v_x, v_y, \omega_z]$ is the commanded velocity.
$\mathbf g \in \mathbb{R}^{\corlGravityDim}$ is the direction of
gravity in the body frame, and
$\mathbf c \in \{0,1\}^{\corlContactDim}$ is the per-foot contact
state. Together these total \corlStateDim{} dimensions.

The network implements a middle-fusion style feature combinator
\cite{feng2021deep}.
\begin{itemize}
  \item \textbf{Robot State Encoder}: The \corlStateDim-dimensional
    state vector is processed by a
    \numberstringnum{\corlStateEncoderLayers}-layer Multi-Layer
    Perceptron (MLP) with a hidden dimension of
    \corlStateEncoderHidden{} and ReLU activations, producing a
    fixed-dimensional latent state representation.
  \item \textbf{Footstep Encoder}: Each candidate's spatial
    coordinates are embedded by a
    \numberstringnum{\corlFootstepEncoderLayers}-layer feedforward
    network with a hidden size of \corlFootstepEncoderHidden{} and
    ReLU activations. The no-action option is not spatial at all; it
    is routed to a separate, learned embedding vector rather than
    through this encoder.
  \item \textbf{Shared Trunk}: The latent state and the candidate
    embedding are concatenated and passed into a
    \numberstringnum{\corlTrunkLayers}-layer shared trunk, with layer
    normalization applied after each linear mapping.
  \item \textbf{Output Heads}: A value head outputs the unnormalized
    logit $\hat{V}$ used for ranking; a duration head outputs
    $t_{\text{swing}}$, scaled by a sigmoid into
    $(\corlSwingDurationMin, \corlSwingDurationMax)$\,s.
\end{itemize}

In the configuration evaluated in this paper, candidates carry no
learned cost input: the footstep evaluation network used in prior
work \cite{Sullivan2025} to pre-score candidates has been removed
entirely, and GaitNet's own value head is the only source of ranking
information. This choice is motivated directly by an ablation in
\cite{Sullivan2025} in which zeroing that cost input \emph{improved}
survival rate over supplying it (see \autoref{sec:results}),
suggesting the pre-computed heuristic
constrained exploration rather than accelerating it.

\subsection{Candidate Generation and Refinement}
\label{subsec:methodology-gaitnet-gradient-ascent}

For each leg, \numberstringnum{\corlCandidatesPerLeg} candidate
positions are drawn from a Gaussian distribution
($\sigma = \corlCandidateStd$\,m) centered on the leg's nominal
footstep offset, then projected onto the nearest valid cell of that
leg's terrain mask via a precomputed Euclidean distance transform.
Because GaitNet's value head is differentiable with respect to a
candidate's spatial coordinates, each of the
\numberstringnum{\corlCandidatesPerLeg} samples is then refined for
\numberstringnum{\corlGradientAscentIters} iterations of gradient
ascent on $\hat V$ at a learning rate of $\corlGradientAscentLR$, with
the candidate re-projected onto the valid mask after every step:
\begin{equation}
  \mathbf{f}_{k+1} = \Pi_{\mathcal{C}} \left( \mathbf{f}_k + \alpha
  \nabla_{\mathbf{f}} \hat{V}(\mathbf{x}_g, \mathbf{f}_k; \theta) \right)
\end{equation}
where $\Pi_{\mathcal{C}}$ projects onto the nearest valid cell of the
mask $\mathcal{C}$, rather than onto a continuous manifold; the mask
resolution is \corlMaskResolutionCm\,cm, so this step is best
understood as a local search over nearby cells rather than continuous
optimization over the workspace. This refinement is a comparatively
minor part of the pipeline: its purpose is to nudge a candidate toward
better-supported footing within its neighborhood, and the terrain
mask, not the gradient step, is what excludes voids and unsteppable
regions from consideration in the first place.

After refinement, GaitNet evaluates the fixed no-action embedding
alongside the
$\corlCandidatesPerLeg \times \corlNumLegs = \corlTotalCandidatesDeploy$
refined candidates (\numberstringnum{\corlCandidatesPerLeg} per leg,
\numberstringnum{\corlNumLegs} legs) and selects the single
highest-valued option. If no-action wins, or if a leg's mask has no
valid candidates, no footstep is issued that tick.

\subsection{Reinforcement Learning Framework}
\label{subsec:methodology-gaitnet-training}

Because identifying optimal footsteps depends on temporal horizons
containing discrete transitions in the contact state $\mathbf{c}$,
generating stable, labeled ground-truth data for supervised
regression is prohibitively difficult. GaitNet is instead trained
with reinforcement learning, using Proximal Policy Optimization (PPO)
\cite{Schulman_proximal_2017}.

The actor network is GaitNet itself. The critic is a smaller,
separately-sized network that scores each candidate through its own
state encoder, footstep encoder, and trunk, then combines the
per-option scores across all sixteen options with a small combining
MLP to produce a single state-value estimate; it does not share
weights with the actor and omits the duration head. Training uses a
fixed candidate set rather than the refinement procedure above: for
each leg, \numberstringnum{\corlTrainCandidatesPerLeg} candidates are
sampled uniformly at random from the valid terrain mask (plus one
no-action option, \corlTrainTotalCandidates{} candidates total), and
the actor's logits over this set define a categorical
distribution. Gradient-based refinement is used only at deployment
time, applied to this same trained value head; training and
deployment therefore draw candidates from different distributions,
which we treat as an intentional design choice rather than an
oversight, since it lets the value head be learned as a general
scorer rather than overfit to the refinement trajectory.

The duration output for the selected candidate is treated as a
Normal distribution with a fixed standard deviation of
\corlDurationExploreStd\,s, concentrating exploration on the discrete
action choice rather than timing. The total log probability of an
action is the sum of the categorical log probability of the selected
candidate and the normal log probability of its sampled duration.
Reward terms balance velocity tracking against stability and a small
per-step action penalty ($\corlActionPenalty$) that discourages
unnecessary stepping; termination on falling or on stepping into
invalid terrain carries a $\corlTerminationPenalty$ penalty. Full
reward and termination tables, command sampling ranges, and PPO
hyperparameters are given in the appendix. The model is trained in
NVIDIA Isaac Lab \cite{mittal_orbit_2023}, with the physics simulation
and policy running on GPU and the per-robot MPC controllers running in
parallel on CPU, for \corlTrainingIterations{} iterations across
\corlParallelEnvs{} parallel environments, with a terrain-difficulty
curriculum that gradually increases the proportion of unsteppable
terrain as training progresses.
